\documentclass{report}
\usepackage{amsmath,amssymb}
\setlength{\parindent}{0mm}
\setlength{\parskip}{1em}
\begin{document}
\begin{center}
\rule{6in}{1pt} \
{\large Geoffrey D. Sanders \\
{\bf Eigenvectors of Matrices Associated with Scale-Free Graphs}}

Center for Applied Scientific Computing \\ Lawrence Livermore National Laboratory \\ 7000 East Avenue \\ Livermore CA 94551
\\
{\tt sanders29@llnl.gov}\end{center}

A scale-free graph is characterized as having vertices with a degree
distribution that follows a power law. Clustering, ranking, and measuring
similarity of vertices on scale-free graphs are all quantities that
network scientists often calculate. We describe a few network
calculations that are obtained by solving linear systems, eigensystems,
nonlinear constrained optimization problems, and tensor decompositions
involving matrices (or tensors) whose sparsity structure is related to an
underlying scale-free graph. In each of these settings, linear solves are
often fundamental calculations performed as subroutines within the larger
iterative solution process. For many large systems of interest, classical
iterative solvers (such as conjugate gradient) converge with
prohibitively slow rates, due to ill-conditioned matrices and small
spectral gap ratios. Scalable preconditioners for this class of problem
are of considerable research interest.

We consider constructing algebraic multigrid (AMG) preconditioners, which
make use of coarse approximation of eigenspaces to accelerate
convergence. We discuss several interesting properties the eigenvectors
of scale-free graphs have that are not present in mesh-like graphs. For
example, these graphs often have a large number of locally supported
eigenvectors (LSEVs), eigenvectors that are each nonzero only in a small,
connected region. Another common related property is that eigenvectors in
a matrix�s near-nullspace are essentially local, having significant
magnitude in a very small region of the graph and rapidly decaying away
from this region. We discuss how such properties can be used to coarsen
graphs with potentially high spectral accuracy and ultimately improve the
scalability of preconditioned iterative methods.

Additionally, these properties are important for understanding how the
quality of a numerical solution relates to the quality of the data mining
task, loosely analogous to the relationship between discretization error
and optimal convergence tolerances in numerical solutions to partial
differential equations.


\end{document}
